\section{Lifting every cross-cup axis indexed by two odd primes}
\label{sec:all-odd-cross-cup-lifts}

Let \(G_0\) act on \(\widehat F(x,y)\) by positive
peripheral automorphisms with common multiplier
\(\lambda\), with \(x\mapsto x^\lambda\),
\(y\mapsto f_g^{-1}y^\lambda f_g\) and \(f_g\in\widehat F(x,y)'\).
Let \(M_0=\ker\lambda\).
For distinct odd primes \(\ell,q\), write
\[
 m=\ell-1,\qquad d=\operatorname{ord}_{\ell}(q),\qquad
 \ell^*=(-1)^{(\ell-1)/2}\ell,\qquad
 q^*=(-1)^{(q-1)/2}q .
\]
Use the cyclotomic quadratic characters \(c_\ell,c_q\)
and the cyclic degree-two classes \(\delta_\ell,\delta_q\).
Write \(h_v(r,s)\) for the additive local Hilbert sign; in
the formulas below \(h_\ell(\ell^*,q^*)=\ell_\ell(q^*)\)
and \(h_q(\ell^*,q^*)=\ell_q(\ell^*)\).

\begin{theorem}\label{thm:all-odd-cross-cup-lifts}
For every distinct pair of odd primes \(\ell,q\), there
is a continuous \(G_0\)-invariant homomorphism
\[
                         \xi_{\ell,q}:M_0\to\mathbf F_2
\]
whose transgression on \(\lambda(G_0)\), pulled back from the
finite cyclotomic quotient, is
\begin{equation}
 e_{c_\ell,c_q}
   =c_\ell c_q+
        h_\ell(\ell^*,q^*)\delta_\ell+
        h_q(\ell^*,q^*)\delta_q.                          \label{eq:all-odd-cross-cup-1}
\end{equation}
Thus every such class has zero inflation to full
profinite GT. The restriction of \(\xi_{\ell,q}\) to
\(G_K\), for \(K=\mathbf Q(\mu_\infty,2^{1/\infty})\),
is exactly the arithmetic basis character
\(\kappa_{\ell,q}\) transgressing to \eqref{eq:all-odd-cross-cup-1}.

In particular all the odd-prime-pair axes of the complete
arithmetic target lie in the image of \(X_2\).
The characters indexed by \(\{a,c_p\}\) are not covered
by this theorem.
\end{theorem}

We give both an explicit beta formula and the complete
finite central extension calculation. The order of the
two primes records which cyclic cover and which
coefficient field are used; no equality between
\(\xi_{\ell,q}\) and \(\xi_{q,\ell}\) on full GT is assumed.

\subsection{The cover and its symplectic representation}

Put
\[
 U_\ell=\ker\bigl(\widehat F(x,y)\to C_\ell,\quad x,y\mapsto1\bigr).
\]
The associated degree-\(\ell\) cover of the oriented
pair of pants has exactly three cusps, since
\(1,1,-2\) are generators of \(C_\ell\).
Its compact genus is \((\ell-1)/2\).
Define \(V\) to be its first homology over \(\mathbf F_q\),
modulo the peripheral homology subspace.
Then
\[
 \dim_{\mathbf F_q}V=m,
\]
and the deck generator \(D\) has each nontrivial
\(\ell\)-th-root character exactly once.

The positive peripheral action gives
\begin{equation}
 \rho_{\ell,q}:G_0\to\operatorname{GSp}(V),\qquad
 \rho_{\ell,q}(g)D\rho_{\ell,q}(g)^{-1}
                  =D^{\lambda_g\bmod\ell},                \label{eq:all-odd-cross-cup-2}
\end{equation}
with similitude multiplier \(\lambda_g\bmod q\).
Here is the algebraic verification of the multiplier,
to include its full scope. The surface presentation is
\[
       \prod_{i=1}^{(\ell-1)/2}[a_i,b_i]\,c_1c_2c_3=1 .
\]
There is one cusp above each original cusp.
A peripheral conjugator in \(\widehat F(x,y)\) can
therefore be moved into \(U_\ell\) by multiplying it by
a suitable power of that original peripheral generator.
Thus each \(c_i\) is sent to a \(U_\ell\)-conjugate of
\(c_i^\lambda\).
In the class-two, exponent-\(q\) quotient, project
commutators to \(\bigwedge^2 V\).
Conjugation changes a peripheral central coordinate
only by a wedge with a peripheral homology class,
which vanishes after projection.
The \(\lambda\)-th power multiplies that projected
coordinate by \(\lambda\bmod q\).
Consequently the symplectic bivector
\(\sum_i a_i\wedge b_i\), and hence its inverse alternating
form, has multiplier \(\lambda\bmod q\).
The deck spectrum follows directly from the chain
complex \(\mathbf F_q[C_\ell]^2\to\mathbf F_q[C_\ell]\):
each nontrivial character has a one-dimensional kernel,
while the two trivial classes are the peripheral subspace.
This proves \eqref{eq:all-odd-cross-cup-2} without an additional covering or filling
equivariance assumption.

Let \(E\) be the normalizer of \(\langle D\rangle\)
in \(\operatorname{GSp}(V)\), and let
\begin{equation}
                 \pi:E\to C_m\times\mathbf F_q^\times
                                                               \label{eq:all-odd-cross-cup-3}
\end{equation}
record the deck exponent and similitude multiplier.
We now describe its kernel \(T\) and an invariant
quadratic character \(\chi:T\to\mathbf F_2\).

\subsection{Even Frobenius order}

Suppose \(d=2s\), and put \(r=m/d\).
The \(r\) Frobenius orbits are self-dual.
Write
\[
 k=\mathbf F_{q^{2s}},\qquad k_0=\mathbf F_{q^s},\qquad
 \bar z=z^{q^s}.
\]
Then \(V=k^r\) as an \(\mathbf F_q\)-space.
Each summand has alternating form
\(\operatorname{Tr}_{k/\mathbf F_q}(d_i z\bar w)\),
where \(\bar d_i=-d_i\).
Rescaling by elements of prescribed \(k/k_0\)-norm
makes all \(d_i\) equal to a common \(d_0\).
Thus
\begin{equation}
 T=\prod_{i=1}^r\ker(N_{k/k_0}:k^\times\to k_0^\times),
 \qquad
 (-1)^{\chi(t)}=\prod_{i=1}^r t_i^{(q^s+1)/2}.            \label{eq:all-odd-cross-cup-4}
\end{equation}
Every factor has even order \(q^s+1\), so \(\chi\ne0\).

Choose a primitive residue \(a\bmod\ell\).
Write \(q=a^{rh}\), with \(\gcd(h,d)=1\), and choose
\(k_1\) modulo \(d\) such that \(a^r=q^{k_1}\).
Then \(\gcd(k_1,d)=1\), so \(k_1\) is odd.
Using orbit representatives \(1,a,\ldots,a^{r-1}\),
a lift of the deck exponent \(a\) cyclically permutes
the \(r\) summands, with a \(q^{k_1}\)-Frobenius twist
on the last passage.

Such a lift can be chosen symplectic.
Indeed, the plain semilinear permutation preserves
the common form on each unwrapped summand.
On the wrapped summand its form coefficient is changed
from \(d_0\) to \(d_0^{q^{-k_1}}\).
Their ratio belongs to \(k_0^\times\).
Multiplying that summand by an element with the
appropriate \(k/k_0\)-norm restores the form, since
the norm map is onto. Denote the resulting symplectic
lift by \(J\).

Choose a generator \(g\) of \(k^\times\), and take
\[
 v=g^{(q^s-1)/(q-1)},\qquad \lambda=Nv.
\]
Then \(\lambda\) generates \(\mathbf F_q^\times\).
The uniform scalar \(L=(v,\ldots,v)\) is a similitude
of multiplier \(\lambda\).
Thus \eqref{eq:all-odd-cross-cup-3} is surjective.

The character \eqref{eq:all-odd-cross-cup-4} is invariant under every permutation
of its factors and every Frobenius automorphism.
Scalar multiplication used in the lift \(J\) commutes
with \(T\). Hence \(\chi\) is \(E\)-invariant.

The scalar correction in \(J\) also commutes with \(L\).
Therefore their commutator is the identity in every
factor except the wrapped one, where it is
\(v^{q^{k_1}-1}\).
This has norm one, and
\begin{equation}
 \begin{split}
 \left(v^{q^{k_1}-1}\right)^{(q^s+1)/2}
 &=g^{\,\frac{q^{2s}-1}{2}
                    \frac{q^{k_1}-1}{q-1}}\\
 &=(-1)^{\,1+q+\cdots+q^{k_1-1}}
 =-1 .
 \end{split}                                               \label{eq:all-odd-cross-cup-5}
\end{equation}
Thus \(\chi([J,L])=1\).
Also
\begin{equation}
                         \chi(L^{q-1})=r\bmod2.            \label{eq:all-odd-cross-cup-6}
\end{equation}

\subsection{Odd Frobenius order}

Suppose \(d\) is odd, and put \(r=m/(2d)\) and
\(k=\mathbf F_{q^d}\).
Frobenius orbits occur in distinct dual pairs.
The representation is a sum of \(r\) hyperbolic
\(k\)-pairs, with form
\[
 \sum_{i=1}^r
 \operatorname{Tr}_{k/\mathbf F_q}(x_i y_i'-y_i x_i').
\]
Rescaling paired coordinates gives this normalization.
The torus and its character are
\begin{equation}
 T=(k^\times)^r,\qquad
 t_i:(x_i,y_i)\mapsto(t_i x_i,t_i^{-1}y_i),\qquad
 (-1)^{\chi(t)}
      =\prod_{i=1}^r t_i^{(q^d-1)/2}.                     \label{eq:all-odd-cross-cup-7}
\end{equation}
Again \(\chi\ne0\).

For a primitive residue \(a\bmod\ell\), the quotient
by \(\langle q,-1\rangle\) has order \(r\), and
\[
                            a^r=-q^{k_1}
\]
for an integer \(k_1\).
Thus the primitive deck multiplier cycles the \(r\)
pairs and on the last passage exchanges the two
dual coordinates with a Frobenius twist.
A sign on one of the two exchanged coordinates
makes this last operation symplectic.
This gives a symplectic lift \(J\).

Take \(L\) on every pair to be
\[
                              (x,y)\mapsto(x,\lambda y),
\]
where \(\lambda\) generates \(\mathbf F_q^\times\).
It is a similitude of multiplier \(\lambda\), proving
surjectivity of \eqref{eq:all-odd-cross-cup-3}.
Permutations, Frobenius twists, and inversion preserve
the sum of quadratic characters in \eqref{eq:all-odd-cross-cup-7}, so \(\chi\)
is invariant under \(E\).

The commutator \([J,L]\) is trivial except on the
wrapped pair, where its torus parameter is
\(\lambda\) or its inverse.
Its quadratic sign in \(k\) is
\begin{equation}
               \lambda^{(q^d-1)/2}=(-1)^d=-1.             \label{eq:all-odd-cross-cup-8}
\end{equation}
Thus the mixed coefficient is again one.
Here \(L^{q-1}=1\), so
\begin{equation}
                              \chi(L^{q-1})=0.             \label{eq:all-odd-cross-cup-9}
\end{equation}

\subsection{The diagonal coefficients and Hilbert corrections}

In both cases \(\ker\chi\) is normal in \(E\), and
\begin{equation}
 1\to C_2\to E/\ker\chi\to C_m\times\mathbf F_q^\times\to1
                                                               \label{eq:all-odd-cross-cup-10}
\end{equation}
is a central extension.
The degree-two cohomology of the base has the three
basis classes \(\delta_\ell,\delta_q,c_\ell c_q\).
Equations \eqref{eq:all-odd-cross-cup-5} and \eqref{eq:all-odd-cross-cup-8} show that the mixed coefficient
is one. Equations \eqref{eq:all-odd-cross-cup-6} and \eqref{eq:all-odd-cross-cup-9} give the coefficient
\(\beta\) of \(\delta_q\):
\begin{equation}
 \beta=
 \begin{cases}
 r\bmod2,&d\text{ even},\\
 0,&d\text{ odd}.
 \end{cases}                                               \label{eq:all-odd-cross-cup-11}
\end{equation}

There is a uniform exact computation of the other
diagonal coefficient. Over an algebraic closure,
\(J\) permutes the \(m\) distinct deck eigenlines
in a single \(m\)-cycle.
Since \(J\) is symplectic, \(\det J=1\).
Since \(m\) is even, the sign of that permutation
is \(-1\). The product of the scalars around the
cycle is therefore \(-1\), and
\begin{equation}
                                J^m=-I.                   \label{eq:all-odd-cross-cup-12}
\end{equation}
In particular one must not replace \(J^m\) by the
identity before applying the central character.
The coefficient \(\alpha\) of \(\delta_\ell\) is
\begin{equation}
 \alpha=\chi(-I)=
 \begin{cases}
 r(q^s+1)/2\bmod2,&d=2s,\\
 r(q^d-1)/2\bmod2,&d\text{ odd}.
 \end{cases}                                               \label{eq:all-odd-cross-cup-13}
\end{equation}
Thus the class of \eqref{eq:all-odd-cross-cup-10} is exactly
\begin{equation}
                     c_\ell c_q+\alpha\delta_\ell+\beta\delta_q.
                                                               \label{eq:all-odd-cross-cup-14}
\end{equation}

Put \(b_\ell=(\ell-1)/2\bmod2\) and
\(b_q=(q-1)/2\bmod2\).
Equations \eqref{eq:all-odd-cross-cup-11}--\eqref{eq:all-odd-cross-cup-13} simplify to
\begin{equation}
 \beta=\frac{1-\left(\frac q\ell\right)}2,\qquad
                         \alpha=\beta+b_\ell b_q.          \label{eq:all-odd-cross-cup-15}
\end{equation}
For the first identity, if \(d\) is odd then \(q\)
is a square modulo \(\ell\).
If \(d=2s\), writing \(q=a^{rh}\) as above gives
\((q/\ell)=(-1)^r\), because \(h\) is odd.
For the second identity, in odd degree
\(b_\ell\equiv r\pmod2\), and
\((q^d-1)/2\equiv b_q\pmod2\).
In even degree, distinguish \(q\equiv1\) or \(3\pmod4\)
and use \(b_\ell\equiv rs\pmod2\).
This gives precisely \eqref{eq:all-odd-cross-cup-15}.

The Hilbert symbol of \((\ell^*,q^*)\) has no finite
ramification outside \(\ell,q\):
both discriminants are odd and one modulo four,
so the symbol at two is also one.
At \(q\), its sign is
\((\ell^*/q)=(q/\ell)\).
At \(\ell\), its sign is
\((q^*/\ell)=(-1)^{b_\ell b_q}(q/\ell)\).
These are the usual odd-prime norm-residue calculations
and quadratic reciprocity
\cite[III, Theorem~4.4, pp.~113--114; Remark~4.7, p.~115]{MilneCFT}.
Thus \(\alpha,\beta\) are exactly the two Hilbert
corrections in \eqref{eq:all-odd-cross-cup-1}.
This identifies the entire extension class directly;
arithmetic Brauer vanishing is not used to infer
vanishing on GT.

\subsection{The resulting characters and their arithmetic restrictions}

Composing \eqref{eq:all-odd-cross-cup-2} with \eqref{eq:all-odd-cross-cup-10} gives a homomorphism from
\(G_0\) lifting its finite cyclotomic quotient
\[
       g\longmapsto
        (\lambda_g\bmod\ell,\lambda_g\bmod q).
\]
Consequently the class \eqref{eq:all-odd-cross-cup-14}, hence \eqref{eq:all-odd-cross-cup-1}, inflates to
zero on \(G_0\).
For \(g\in M_0\), \(\rho_{\ell,q}(g)\in T\), and we set
\[
                       \xi_{\ell,q}(g)=
                                  \chi(\rho_{\ell,q}(g)).
\]
This is continuous, additive, and \(G_0\)-invariant.

Here are explicit formulas in the normalized beta
convention of \cite[\S2.3, p.~101]{Furusho2003}.
Choose a primitive \(\ell\)-th root \(\zeta\) in
\(\mathbf F_{q^d}\).
If \(d=2s\), choose representatives \(a_i\) for
\((\mathbf Z/\ell)^\times/\langle q\rangle\); then
\begin{equation}
 (-1)^{\xi_{\ell,q}(g)}
   =\prod_{i=1}^{r}
       B_g(\zeta^{a_i},\zeta^{a_i})^{(q^s+1)/2}.            \label{eq:all-odd-cross-cup-16}
\end{equation}
Each factor before exponentiation has norm one from
\(\mathbf F_{q^{2s}}\) to \(\mathbf F_{q^s}\).
If \(d\) is odd, choose representatives for
\((\mathbf Z/\ell)^\times/\langle q,-1\rangle\); then
\begin{equation}
 (-1)^{\xi_{\ell,q}(g)}
   =\prod_{i=1}^{r}
       B_g(\zeta^{a_i},\zeta^{a_i})^{(q^d-1)/2}.            \label{eq:all-odd-cross-cup-17}
\end{equation}
In this case the reciprocal deck character has the
inverse scalar.
Changing representatives or the primitive root only
permutes factors and applies Frobenius or inversion,
which preserve the displayed signs.
The commutator class generates each nontrivial
character line of the cyclic cover, so these scalars
are exactly the homology scalars used above.

Naturality of transgression gives \eqref{eq:all-odd-cross-cup-1} on the quotient
\(A=G/N\).
The arithmetic transgression on \(G_K\) is injective,
and its basis character \(\kappa_{\ell,q}\) has that
same class. Hence
\[
                       \xi_{\ell,q}|_{G_K}=\kappa_{\ell,q}.
\]
This completes the proof.

\begin{remark}[Comparing the two ordered constructions]
The characters \(\xi_{\ell,q}|_N\) and
\(\xi_{q,\ell}|_N\) have the same arithmetic restriction,
because their transgression is the symmetric class
\(e_{c_\ell,c_q}\).
Therefore
\[
          \varepsilon_{\ell,q}
             =\left.(\xi_{\ell,q}+\xi_{q,\ell})\right|_N
\]
is \(G\)-invariant and vanishes on \(N\cap H\).
The established factorization \(G=NH\) consequently
gives a unique continuous character \(G\to\mathbf F_2\)
which is \(\varepsilon_{\ell,q}\) on \(N\) and zero on
\(H\): define its value at \(nh\) to be
\(\varepsilon_{\ell,q}(n)\).
Agreement on the intersection makes this well-defined,
and invariance makes it a homomorphism.
This constructs explicit relative characters but does
not prove that any of them is nonzero.
\end{remark}
